% Overleaf-Datei (Name zu bestaetigen), Spiegel in docs/overleaf/.
% Nur zusammen mit der Overleaf-Datei aendern.
% Inhalt: Abstract.

\section*{Abstract}
\addcontentsline{toc}{section}{Abstract}

Machine-learned interatomic potentials for reactions are trained on DFT
labels, and the level of theory of those labels bounds their accuracy.
The common way to raise the level of theory is relabelling: the geometries of an
existing dataset are kept and only their energies and forces are
recomputed at a higher level. This thesis asks two questions about
that practice: whether a model can be lifted to a higher level by
relabelling only a subset of its training data (RQ1), and whether a
model can learn a surface from paths that were never relaxed on that
surface (RQ2).

For RQ1, a MACE model is trained on Transition1x
(\mbox{$\omega$B97X/6-31G(d)}), and a correction head on top of it
is trained to predict the difference to \mbox{$\omega$B97M-V/def2-TZVP}
from about 80\,000 geometries, under 1\,\% of the dataset. It is tested on 30 reactions grouped by low, mid and high
multireference character. The force error drops to about a third on all reactions, and the
corrected model locates the transition states accurately. Barriers improve
threefold for reactions of low multireference character; for reactions
of mid and high multireference character the corrected barriers
overshoot, mostly through the error of the MACE model on its own
level, which the head cannot correct. The answer to RQ1 is yes, with the reservation of
the reactions of high multireference character.

In RQ1 the labels were raised to a higher level of theory while the
spin formalism stayed restricted. In RQ2 we ask whether relabelling still
works when the spin formalism changes as well. For that we turn to
models trained on OMol25, the current frontier of general-purpose
molecular potentials. OMol25 gave unrestricted labels at
\mbox{$\omega$B97M-V/def2-TZVPD} to reaction paths that had been
generated restricted, the Transition1x paths among them.

Published
benchmarks of these models report high success rates in
transition-state search. None has tested the models systematically at
broken-symmetry transition states. This work does so for three
of the models: UMA-S, UMA-M and eSEN. We let them search the
transition states of 45 Transition1x reactions. Each transition state
the models deliver is evaluated with an unrestricted
DFT single point and a stability analysis, which sorts the transition
states into closed-shell and broken-symmetry. The models predict single-point energies well for both groups. But
their force errors are two to three times larger at the
broken-symmetry transition states. The structures they deliver there
are therefore not stationary points of the unrestricted surface, with
residual forces up to 3.4 times those of the closed-shell group. This
inability to locate the transition state makes the barriers at the
broken-symmetry reactions more than ten times less reliable than at the
closed-shell ones. Two mechanisms are candidates for these errors. The first is the training data: the models learned from transition
states found on the restricted surface, and these are furthest from
being stationary points of the unrestricted surface at the
broken-symmetry reactions. The second is the shape of the unrestricted
surface, whose curvature changes abruptly where the broken-symmetry
solution sets in, which may make it hard for a smooth model to learn.
The data do not decide between the two. We also find that the
prediction quality of MACE and of the OMol25 models degrades with
multireference character independently of broken symmetry.

For practice, a stability analysis should precede any further use of a
transition state whose multireference character cannot be ruled out.
Fine-tuning the models on paths relaxed on the unrestricted surface
would separate the two mechanisms and is left to future work.
